\documentclass[12pt,a4paper,oneside,pdftex]{report}
\usepackage[latin1]{inputenc}
\usepackage{lmodern}
\usepackage[OT1]{fontenc}
\usepackage[finnish,swedish,english]{babel}


\usepackage[square,sort&compress,numbers]{natbib}
\usepackage{eurosym} 

\usepackage{verbatim}

\usepackage{longtable}

\usepackage{subfigure}

\usepackage[medium]{titlesec}

\usepackage{tikz}
% You also need to specify which TikZ libraries you use
\usetikzlibrary{positioning}
\usetikzlibrary{calc}
\usetikzlibrary{arrows}
\usetikzlibrary{decorations.pathmorphing,decorations.markings}
\usetikzlibrary{shapes}
\usetikzlibrary{patterns}


% The aalto-thesis package provides typesetting instructions for the
% standard master's thesis parts (abstracts, front page, and so on)
% Load this package second-to-last, just before the hyperref package.
% Options that you can use: 
%   mydraft - renders the thesis in draft mode. 
%             Do not use for the final version. 
%   doublenumbering - [optional] number the first pages of the thesis
%                     with roman numerals (i, ii, iii, ...); and start
%                     arabic numbering (1, 2, 3, ...) only on the 
%                     first page of the first chapter
%   twoinstructors  - changes the title of instructors to plural form
%   twosupervisors  - changes the title of supervisors to plural form
\usepackage[mydraft]{aalto-thesis}
%\usepackage[mydraft,doublenumbering]{aalto-thesis}
%\usepackage{aalto-thesis}


\RequirePackage[pdftex]{hyperref}
\hypersetup{colorlinks=false,raiselinks=false,breaklinks=true}
\hypersetup{pdfborder={0 0 0}}
\hypersetup{bookmarksnumbered=true}
\hypersetup{bookmarksopen=true,bookmarksopenlevel=1}

\newcommand{\TITLE}{Lean software development and the effects of working with unfamiliar technology}
\newcommand{\FTITLE}{}
\newcommand{\STITLE}{}
\newcommand{\SUBTITLE}{A case study}
\newcommand{\FSUBTITLE}{}
\newcommand{\SSUBTITLE}{}
\newcommand{\DATE}{December 1, 2015}
\newcommand{\FDATE}{1. joulukuuta 2015}
\newcommand{\SDATE}{Den 1 December 2015}

% Supervisors and instructors
% ------------------------------------------------------------------
% If you have two supervisors, write both names here, separate them with a 
% double-backslash (see below for an example)
% Also remember to add the package option ``twosupervisors'' or
% ``twoinstructors'' to the aalto-thesis package so that the titles are in
% plural.
% Example of one supervisor:
%\newcommand{\SUPERVISOR}{Professor Antti Yl�-J��ski}
%\newcommand{\FSUPERVISOR}{Professori Antti Yl�-J��ski}
%\newcommand{\SSUPERVISOR}{Professor Antti Yl�-J��ski}
% Example of twosupervisors:
\newcommand{\SUPERVISOR}{Professor Marjo Kauppinen}
\newcommand{\FSUPERVISOR}{Professori Marjo Kauppinen}
\newcommand{\SSUPERVISOR}{Professor Marjo Kauppinen}

% If you have only one instructor, just write one name here
\newcommand{\INSTRUCTOR}{Eeva \\
Suvi}
\newcommand{\FINSTRUCTOR}{Eeva \\
Suvi}
\newcommand{\SINSTRUCTOR}{Eeva \\
Suvi}
% If you have two instructors, separate them with \\ to create linefeeds
% \newcommand{\INSTRUCTOR}{Olli Ohjaaja M.Sc. (Tech.)\\
%  Elli Opas M.Sc. (Tech)}
%\newcommand{\FINSTRUCTOR}{Diplomi-insin��ri Olli Ohjaaja\\
%  Diplomi-insin��ri Elli Opas}
%\newcommand{\SINSTRUCTOR}{Diplomingenj�r Olli Ohjaaja\\
%  Diplomingenj�r Elli Opas}

% If you have two supervisors, it is common to write the schools
% of the supervisors in the cover page. If the following command is defined,
% then the supervisor names shown here are printed in the cover page. Otherwise,
% the supervisor names defined above are used.
\newcommand{\COVERSUPERVISOR}{Professor Marjo Kauppinen, Aalto University}

% The same option is for the instructors, if you have multiple instructors.
% \newcommand{\COVERINSTRUCTOR}{Olli Ohjaaja M.Sc. (Tech.), Aalto University\\
%  Elli Opas M.Sc. (Tech), Aalto SCI}


% Other stuff
% ------------------------------------------------------------------
\newcommand{\PROFESSORSHIP}{Software Engineering and Business}
\newcommand{\FPROFESSORSHIP}{Ohjelmistotuotanto ja -liiketoiminta}
\newcommand{\SPROFESSORSHIP}{Programvaruproduktion och -aff�rsverksamhet}
% Professorship code is the same in all languages
\newcommand{\PROFCODE}{T-76}
\newcommand{\KEYWORDS}{innovation, services, service design}
\newcommand{\FKEYWORDS}{innovaatio, palvelut, palvelusuunnittelu}
\newcommand{\SKEYWORDS}{innovation, tj�nster, service design}
\newcommand{\LANGUAGE}{English}
\newcommand{\FLANGUAGE}{Englanti}
\newcommand{\SLANGUAGE}{Engelska}

% Author is the same for all languages
\newcommand{\AUTHOR}{Oskar Ehnstr�m}


% Currently the English versions are used for the PDF file metadata
% Set the PDF title
\hypersetup{pdftitle={\TITLE\ \SUBTITLE}}
% Set the PDF author
\hypersetup{pdfauthor={\AUTHOR}}
% Set the PDF keywords
\hypersetup{pdfkeywords={\KEYWORDS}}
% Set the PDF subject
\hypersetup{pdfsubject={Master's Thesis}}


% Layout settings
% ------------------------------------------------------------------

% Bibliography style
% acm style gives you a basic reference style. It works only with numbered
% references.
\bibliographystyle{acm}
% Plainnat is a plain style that works with both numbered and name citations.
% \bibliographystyle{plainnat}


% Extra hyphenation settings
\hyphenation{di-gi-taa-li-sta yksi-suun-tai-sta}



% The preamble ends here, and the document begins. 
% Place all formatting commands and such before this line.
% ------------------------------------------------------------------
\begin{document}
% This command adds a PDF bookmark to the cover page. You may leave
% it out if you don't like it...
\pdfbookmark[0]{Cover page}{bookmark.0.cover}
% This command is defined in aalto-thesis.sty. It controls the page 
% numbering based on whether the doublenumbering option is specified
\startcoverpage

% Cover page
% ------------------------------------------------------------------
% Options: finnish, english, and swedish
% These control in which language the cover-page information is shown
\coverpage{english}


% Abstracts
% ------------------------------------------------------------------
% Include an abstract in the language that the thesis is written in,
% and if your native language is Finnish or Swedish, one in that language.

% Abstract in English
% ------------------------------------------------------------------
\thesisabstract{english}{
A dissertation or thesis is a document submitted in support of candidature
for a degree or professional qualification presenting the author's research and
findings. In some countries/universities, the word thesis or a cognate is used
as part of a bachelor's or master's course, while dissertation is normally
applied to a doctorate, whilst, in others, the reverse is true.

\fixme{Abstract text goes here (and this is an example how to use fixme).} 
Fixme is a command that helps you identify parts of your thesis that still
require some work. When compiled in the custom \texttt{mydraft} mode, text
parts tagged with fixmes are shown in bold and with fixme tags around them. When
compiled in normal mode, the fixme-tagged text is shown normally (without
special formatting). The draft mode also causes the ``Draft'' text to appear on
the front page, alongside with the document compilation date. The custom
\texttt{mydraft} mode is selected by the \texttt{mydraft} option given for the
package \texttt{aalto-thesis}, near the top of the \texttt{thesis-example.tex}
file.

The thesis example file (\texttt{thesis-example.tex}), all the chapter content
files (\texttt{1introduction.tex} and so on), and the Aalto style file
(\texttt{aalto-thesis.sty}) are commented with explanations on how the Aalto
thesis works. The files also contain some examples on how to customize various
details of the thesis layout, and of course the example text works as an
example in itself. Please read the comments and the example text; that should
get you well on your way!}

% Abstract in Finnish
% ------------------------------------------------------------------
\thesisabstract{finnish}{
  Lorem ipsum Tempor tempor labore Ut reprehenderit tempor irure incididunt non labore irure dolore consectetur esse sit magna culpa ad consequat sit Ut velit veniam dolore fugiat sed nostrud reprehenderit.

  Lorem ipsum Magna non sint incididunt laboris Ut proident exercitation dolore non eu adipisicing ullamco occaecat cupidatat Duis labore eiusmod veniam nisi dolor ea est deserunt exercitation aliquip commodo anim magna dolore anim ad mollit enim officia proident nisi quis.

  Lorem ipsum Culpa enim sed tempor velit incididunt dolor aliqua consectetur ut quis officia consectetur proident magna sunt nulla cupidatat laboris et est nisi dolore exercitation consectetur est consequat ea aliqua officia Excepteur reprehenderit aliqua cillum.
  }

% Abstract in Swedish
% ------------------------------------------------------------------
\thesisabstract{swedish}{
  Lorem ipsum Tempor tempor labore Ut reprehenderit tempor irure incididunt non labore irure dolore consectetur esse sit magna culpa ad consequat sit Ut velit veniam dolore fugiat sed nostrud reprehenderit.

  Lorem ipsum Magna non sint incididunt laboris Ut proident exercitation dolore non eu adipisicing ullamco occaecat cupidatat Duis labore eiusmod veniam nisi dolor ea est deserunt exercitation aliquip commodo anim magna dolore anim ad mollit enim officia proident nisi quis.

  Lorem ipsum Culpa enim sed tempor velit incididunt dolor aliqua consectetur ut quis officia consectetur proident magna sunt nulla cupidatat laboris et est nisi dolore exercitation consectetur est consequat ea aliqua officia Excepteur reprehenderit aliqua cillum.
}


% Acknowledgements
% ------------------------------------------------------------------
% Select the language you use in your acknowledgements
\selectlanguage{english}

% Uncomment this line if you wish acknoledgements to appear in the 
% table of contents
%\addcontentsline{toc}{chapter}{Acknowledgements}

% The star means that the chapter isn't numbered and does not 
% show up in the TOC
\chapter*{Acknowledgements}

TODO: Thank people here

\vskip 10mm

\noindent Espoo, \DATE
\vskip 5mm
\noindent\AUTHOR

% Acronyms
% ------------------------------------------------------------------
% Use \cleardoublepage so that IF two-sided printing is used 
% (which is not often for masters theses), then the pages will still
% start correctly on the right-hand side.
%\cleardoublepage
% Example acronyms are placed in a separate file, acronyms.tex
%\addcontentsline{toc}{chapter}{Abbreviations and Acronyms}
\chapter*{Abbreviations and Acronyms}

% The longtable environment should break the table properly to multiple pages, 
% if needed

\noindent
\begin{longtable}{@{}p{0.25\textwidth}p{0.7\textwidth}@{}}
2k/4k/8k mode & COFDM operation modes \\
3GPP & 3rd Generation Partnership Project \\ 
ESP & Encapsulating Security Payload; An IPsec security protocol \\ 
FLUTE  & The File Delivery over Unidirectional Transport protocol \\ 
e.g.& for example (do not list here this kind of common acronymbs or abbreviations, but only those that are essential for understanding the content of your thesis. \\ 
note & Note also, that this list is not compulsory, and should be omitted if you have only few abbreviations

\end{longtable}


% Table of contents
% ------------------------------------------------------------------
\cleardoublepage
% This command adds a PDF bookmark that links to the contents.
% You can use \addcontentsline{} as well, but that also adds contents
% entry to the table of contents, which is kind of redundant.
% The text ``Contents'' is shown in the PDF bookmark. 
\pdfbookmark[0]{Contents}{bookmark.0.contents}
\tableofcontents

% List of tables
% ------------------------------------------------------------------
% You only need a list of tables for your thesis if you have very 
% many tables. If you do, uncomment the following two lines.
% \cleardoublepage
% \listoftables

% Table of figures
% ------------------------------------------------------------------
% You only need a list of figures for your thesis if you have very 
% many figures. If you do, uncomment the following two lines.
% \cleardoublepage
% \listoffigures

% The following label is used for counting the prelude pages
\label{pages-prelude}
\cleardoublepage

%%%%%%%%%%%%%%%%% The main content starts here %%%%%%%%%%%%%%%%%%%%%
% ------------------------------------------------------------------
% This command is defined in aalto-thesis.sty. It controls the page 
% numbering based on whether the doublenumbering option is specified
\startfirstchapter

% Add headings to pages (the chapter title is shown)
\pagestyle{headings}

% The contents of the thesis are separated to their own files.
% Edit the content in these files, rename them as necessary.
% ------------------------------------------------------------------
 % -*- root: thesis.tex -*-
\chapter{Introduction}
\label{chapter:intro}

\section{Background}
\label{section:background}

Futurice is a software consult agency developing digital solutions for their clients. Their methodology has alwasy been agile and has long focused on the overall experience of the end-users. They aim to take their customers through the whole lifecycle of a product, from ideation to lifecycle management. The process Futurice uses for this holistic approach is called Lean Service Creation (LSC).


\section{Research problem and question}
\label{section:problem}

The research problem is defined as follows:\\

\textit{How does a new service creation project that uses new and untested technology differ from a traditional service creation project that uses familiar technologies and concepts?}\\

To investigare this problem three research questions have been set up in table \ref{tbl:questions}.


\begin{table}
  \begin{tabular}{p{200pt} | p{70pt} | p{70pt}}
    Question & Litterature review & Empirical study \\
    \hline
    What are the currently available frameworks for new service creation? & x & \\
    What are the success factors of a new service creation project involving novel technology? &  & x \\
    What are the challenges when executing a new service creation project using novel technology? &  & x \\
  \end{tabular}
  \caption{Research questions and their respective sections}
  \label{tbl:questions}
\end{table}

\section{Scope}
\label{section:scope}




\section{Structure of the Thesis}
\label{section:structure}

The thesis is split into two distinct sections. First, chapter \ref{chapter:litterature} covers service design and innovation in general. Second, chapter \ref{chapter:empirical} covers the findings of the innovation culture at Futurice and how exposure to new technology has affected their new service creation work.




 % -*- root: thesis.tex -*-

\chapter{Litterature review}
\label{chapter:litterature}

In this section I will present the existing litterature and what general guidelines or conclusions it presents. I will refer to the main concepts presented in books and compare those findings to case studies that have implemented the relevant concepts in order to find some commonalities if there are any.

I will also present the process for finding and analyzing the material I have chosen.

\begin{itemize}
  \item{How was the litterature review conducted?}
  \begin{itemize}
    \item{Scholar}
    \item{Scholar through related articles of known good articles}
    \item{Key articles then their sources}
    \item{Same author, look for more}
    \item{Same conference/journal, look for more}
  \end{itemize}
  \
\end{itemize}

Keywords: \textit{new service creation, digital service creation, service-dominant, design thinking, new service development}

\section{The origin of lean thinking}
\label{section:leanorigins}


Even though the traditional wisdom is that the Japanese car manufacturers had a significant advantage over western competitors due to the lean methodologies they implemented there is some controcersy over wether this was in fact the case. Dyb� \& Sharp argue that by examining the facts and taking automation into account the Japanese did not have a superior organizational advantage. \cite{Dyba2012WhatS}

\section{Lean versus Agile}
\label{section:leanvsagile}

\fixme{
Lean should be thought of a set of principles rather than practices. This article has some excellent points and trends to talk about.\cite{Poppendieck2012Lean}
}

End with something about everything being services nowdays and lean/agile being developer focused. The continue to...

\section{Lean Service Creation}
\label{section:lsc}

\fixme{Lean Service Creation is about build measure learn and create somehing new. Tie this into the problem and focus on the new part. Using news devices maybe?}

Lean Service Creation is based on the ideas of ``The Lean Startup'' as described by Eric Reis in his 2011 book.\cite{ries2011lean}

\fixme{How do I get a source for this? Interview?}

\section{Comparison of existing litterature}

In this section I will compare articles regarding case studies focusing on lean software development/management.

Two case studies. This is an early attempt to test lean software development. The article covers an experiment set up to test the validity of using lean principles in software development. \cite{Middleton2001Lean}

Timberline Inc. case study. Probably the first case of adopting lean principles to software development.\cite{Middleton2005Lean}

BBC Worldwide case study. The gist of it is that the performance of the team improved when adopting lean practices, but there were some challenges in fitting the the lean principles with the rest of the company.\cite{Middleton2012Lean}


\section{Research problem and question}
\label{section:problem}

This chapter will end with the research problem and questions.

The research problem is defined as follows:\\

\textit{What are the commonly accepted and used lean software development practices and how do they change when working with new and unfamiliar technology?}\\

To investigare this problem three research questions have been set up in table \ref{tbl:questions}.


\begin{table}
  \begin{tabular}{p{200pt} | p{70pt} | p{70pt}}
    Question & Litterature review & Empirical study \\
    \hline
    What are the currently available best practices for lean software projects? & x & \\
    Which of these best practices need to be adapted when working with new technology? &  & x \\
    How do these best practices need to adapted? &  & x \\
    Which best practices remain valid? &  & x \\
  \end{tabular}
  \caption{Research questions and their respective sections}
  \label{tbl:questions}
\end{table}

 % -*- root: thesis.tex -*-

\chapter{Methods}
\label{chapter:mehods}

In this section I will describe how I conducted the empirical study.

\begin{itemize}
  \item{Interviews}
  \begin{itemize}
    \item{2-3 project members}
    \item{2-3 customers}
    \item{2-3 end-users}
    \item{1-1.5 hours each (5-10min for end-users)}
  \end{itemize}
  \item{Semi-structured interviews}
\end{itemize}



 % -*- root: thesis.tex -*-

\chapter{Results}
\label{chapter:empirical}

In this section I will present the results of the empirical study. This section contains the raw data of the conducted interviews. This section does not analyze or compare the results with the litterature.


 % -*- root: thesis.tex -*-

\chapter{Discussion}
\label{chapter:discussion}

At this point, you will have some insightful thoughts on your implementation
and you may have ideas on what could be done in the future. 
This chapter is a good place to discuss your thesis as a whole and to show your
professor that you have really understood some non-trivial aspects of the
methods you used\ldots



 % -*- root: thesis.tex -*-

\chapter{Conclusions}
\label{chapter:conclusions}

Here I mention the most important findings of the discussion section and the litterary review section.

I also point out how this research can be used in the future and what it's limitations are. (e.g. only one case study)

2 pages 



% Load the bibliographic references
% ------------------------------------------------------------------
% You can use several .bib files:
% \bibliography{thesis_sources,ietf_sources}
\bibliography{sources,Remote}


% Appendices go here
% ------------------------------------------------------------------
% If you do not have appendices, comment out the following lines
\appendix
 % -*- root: thesis.tex -*-

% \chapter{First appendix}
% \label{chapter:first-appendix}

% This is the first appendix. You could put some test images or verbose data in an
% appendix, if there is too much data to fit in the actual text nicely.

% For now, the Aalto logo variants are shown in Figure~\ref{fig:aaltologo}.

% \begin{figure}
% \begin{center}
% \subfigure[In English]{\includegraphics[width=.8\textwidth]{images/aalto-logo-en}}
% \subfigure[Suomeksi]{\includegraphics[width=.8\textwidth]{images/aalto-logo-fi}}
% \subfigure[På svenska]{\includegraphics[width=.8\textwidth]{images/aalto-logo-se}}
% \caption{Aalto logo variants}
% \label{fig:aaltologo}
% \end{center}
% \end{figure}

\chapter{Interview questions}
\label{chapter:interview-questions}

Here goes the questions from the interviews.


% End of document!
% ------------------------------------------------------------------
% The LastPage package automatically places a label on the last page.
% That works better than placing a label here manually, because the
% label might not go to the actual last page, if LaTeX needs to place
% floats (that is, figures, tables, and such) to the end of the 
% document.
\end{document}
